\chapter{两万一千三百年}
\chapterkana{二万一千三百年}
\epigraphbox{“高岸为谷，深谷为陵。”}{——《诗经·小雅·十月之交》}

\kaishou{一根岩芯，两万一千三百年前的一层泥。\par
层下是一种树，层上是另一种。\par
中间没有过渡，像有人把两页纸对齐了粘在一起。}

\section{岩芯}

贝加尔湖底常年沉积着一层一层的湖泥，年复一年，连成一部未曾中断的天然年历。一九九六年，俄美联合钻探队在湖心取得岩芯，编号 BDP-96。\footnote{贝加尔湖钻探计划 BDP-96 孔，一九九六年夏于 Academician Ridge 施工，取芯进尺二百米，此处用其上段。}取芯时用钢管边转边压，将湖底连泥带水一并提起，到船上再剖开管壁，将泥柱按厘米分段描述与分析。如此取出的泥柱不会翻搅，原始层序得以完整保存，这是它能充当年历的前提。

至于年层如何计数，全靠湖泥的季节性差异：春季硅藻繁盛，形成一浅色薄层；秋季陆源碎屑增多，形成一深色薄层，一浅一深便计作一年。这套数法在浅水湖泊往往不可靠，贝加尔湖却行得通，只因湖水极深、底部缺氧且扰动极少。为求稳妥，同一根岩芯上还做了放射性同位素测年，两种方法的结果相差不远，均落在允许的容差之内。\footnote{纹层计数与放射性同位素测年互为校核，是湖相沉积年代学的常规做法；两法若差得远，这一层便不能当年历用。此处两法相差不到三百年，已在本书容差之内。}

到了六点零三米处，这一层的孢粉谱忽然剧变：原本繁茂的针叶林退走，草原植被取而代之，硅藻数量跌落至原先的几分之一，随后才缓缓回升。\footnote{孢粉谱按云杉、松、蒿、藜等几类分别计数，每样不少于三百粒，才够计算百分比。样本密度为每十厘米一个，六点零三米上下各取了三个样。}这一转折层被定在距今约两万一千三百年。

\begin{figure}[htbp]
\centering
\begin{tikzpicture}[scale=1.0, font=\footnotesize]
  \useasboundingbox (-1.4,-4.4) rectangle (9.2,3.4);
  % 深度轴
  \draw[line width=0.4pt] (0,-4.0) -- (0,3.0);
  \foreach \y/\d in {3.0/5.70, 0.12/6.03, -4.0/6.30}{
    \draw[line width=0.3pt] (-0.12,\y) -- (0,\y);
    \node[anchor=east] at (-0.16,\y) {\d};
  }
  \node[anchor=north, rotate=90] at (-0.95,-0.5) {深度（米）};
  % 岩芯柱
  \fill[black!5] (0.35,-4.0) rectangle (1.35,3.0);
  \draw[line width=0.5pt] (0.35,-4.0) rectangle (1.35,3.0);
  \fill[black!60] (0.35,0.05) rectangle (1.35,0.19);
  % 孢粉曲线：云杉（左）、蒿（右）
  \draw[line width=0.5pt] (5.0,-4.0) -- (5.0,3.0);
  \draw[line width=0.5pt] (8.0,-4.0) -- (8.0,3.0);
  \node[anchor=south] at (5.0,3.0) {云杉};
  \node[anchor=south] at (8.0,3.0) {蒿};
  \draw[line width=0.8pt]
    (7.7,-4.0) -- (7.6,-3.0) -- (7.75,-2.0) -- (7.6,-1.0)
    -- (7.7,-0.25) -- (5.3,0.02) -- (5.2,1.0) -- (5.35,2.0) -- (5.2,3.0);
  \draw[line width=0.8pt, dash pattern=on 2pt off 1.6pt]
    (5.3,-4.0) -- (5.25,-3.0) -- (5.4,-2.0) -- (5.25,-1.0)
    -- (5.35,-0.25) -- (7.6,0.02) -- (7.7,1.0) -- (7.55,2.0) -- (7.7,3.0);
  \draw[line width=0.25pt, black!45] (1.35,0.12) -- (4.9,0.12);
\end{tikzpicture}
\caption{六点零三米上下的孢粉曲线。实线为云杉，虚线为蒿；两者在约三厘米的区间内交换位置，图上以一条水平细线标出。}
\figsource{据 BDP-96 孔上段孢粉分析数据重绘。样本间距十厘米，曲线为相邻样点的直线连接，不代表实际的分辨率。}
\end{figure}

单从地质学上看，这不过是一次区域性的植被退缩，与同一根岩芯里其余几处冷暖交替并无本质区别，转折也都在几百年间完成。真正值得留意的，是这一层格外薄。上下各半厘米之内，云杉花粉与蒿属花粉便完成了此消彼长，中间几乎看不见过渡带。按湖底每年的沉积速率折算，半厘米不过相当于几十年，也就是说，那一带的针叶林在短短几十年内退走，此后便再未恢复。

若是气候自身缓慢摆动，森林应当退一步又进半步，在孢粉谱上留下一段来回拉锯的过渡带：云杉比例降到某一水平后略作停顿，回升少许，再继续下降。同一根岩芯中的其他几处气候转折都带有这样的拉锯痕迹，唯独六点零三米这一层没有，仿佛从一种植被谱直接切换成了另一种。然而这一层泥所能说明的也就到此为止：某件事在距今约两万一千三百年发生，来得极快，且波及贝加尔湖一侧；至于那究竟是什么事，泥层里一个字也没有写。

\section{题头}

岩芯只给出了年代，要把年代落到人身上，还得换一批材料：那批在城中发现的金属片，共四百一十二枚，单面刻字，按内容编作六卷，编目时统称《章程总目》。第三章将专门讨论这批材料，这里先取第一枚来看。

第一枚置于卷首，习惯上称作题头。它比其余各枚短得多，仅有三行汉字，刻痕等大，全无涂改。

\begin{archive}[内部参考]{档案 № ФВ-1965-09}
题头一枚，三行。上行：「计口三十一万二千。」中行：「完工之日，十五万。」下行：「存者如右。」
\end{archive}

文字依旧式直行、从右往左读，因此「如右」指的并非上文，而是紧随其后另起的一栏，只是那一栏空着未填。第一行「计口」即核算人口，「完工之日」是指城池落成、居民迁入的那一天。三十一万二千这个数并非虚数，可以从城的形制反推出来：几重环垣、每环宽度、每平方公里容纳多少人，算下来正与这一规模相合。\footnote{城形制与人口的详细推算，见《月之都——〈广寒宫残简〉辑佚与考释》论城形制的一章。该书算的是静海七号，环制与背面本体相同，可以移用。}

第二行则给出了另一个数字：十五万。两数之间整整差出一半有余。多出来的那一半既非死于非命，也非账目错算，而是未能登船：城池最初按三十一万二千人的容量开工，等到完工时，船位却只有十五万个。究竟谁上谁留，题头上只留了「存者如右」四字，后面那一栏至今空白。那片空白究竟是当初刻成时便未及填写，还是后来被人磨去，红外成像下看不出任何打磨痕迹，两种可能只好一并留着。

「船位」这个词很值得停下来细想。人在那座城里定居之前，是乘船上去的；再往前推，他们原本就住在地面上。这层意思题头自己并未明说，是「完工」与「船位」二词并置一处才显露出来的。至于表中的第三个数，即撤离那年在册的十一万九千人，见于名籍卷末枚，将在第四章详细讨论。

\begin{table}[htbp]
\centering
\caption{题头与名籍卷所记的三个数}
\begin{tabular}{@{}llll@{}}
\toprule
数 & 出处 & 是什么 & 与上一个数的差\\
\midrule
312{,}000 & 题头第一行 & 设计容量 & — \\
150{,}000 & 题头第二行 & 完工登船人数 & 162{,}000 \\
119{,}000 & 名籍卷末枚 & 撤离那年在册 & 31{,}000 \\
\bottomrule
\end{tabular}
\tabsource{三个数分见题头一件、名籍卷末枚一件。差数为笔者所减，两份材料本身没有这一栏。}
\end{table}

把三个数摆在一处看，三十一万二千是设计容量，十五万是完工登船人数，十一万九千是两万年后的在册余数，三者之间的差额便是这座城两万年间经历的全部减法。这枚题头还有一处紧要之处，就是它不记年份。城中历法与地面的纪年本不在同一套系统里，换算须经过好几道步骤（第七章会详述），中间稍有偏差便足以吃掉数百年光阴。因此本章只能将两件事并排摆着：岩芯定年约在两万一千三百年，题头所记的完工大约在其前后，至于两者之间究竟相隔数十年还是三百年，目下无从细断。

\section{遗迹}

既从「船位」一词推出他们原居地面，那么在离开之前，地上理应留有工程遗迹。这类遗迹确实不止一处，先看一九七二年测绘过的那一处，它埋在所谓有穷氏旧地的地下。\footnote{「有穷氏旧地」是基金会的编目用词。有穷氏是传世文献中的族名，《左传》襄公四年记后羿「因夏民以代夏政」，指的便是这一族；至于其居地是否确在此处，文献只言在河济之间，未出具体方位。本章借用此名，仅指代该处地下设施，不表示已证其地望。}

判定这批地下设施与天上那座城同源，依靠的是三条旁证。其一在形制：测绘记录报出的尺寸是十四米、六十米、二十米三个整数，这在实测中很少见；若换算作六丈、二十六丈与八丈五，取整之后恰好便是这三个数。\footnote{此处借用秦汉尺度，只因第一章折算静海七号的里与丈已用此套，两处若要对比，必须先统一至同一把尺上。按十尺为丈、一尺二十三点一厘米计，六丈合十三点九米（取整十四米），二十六丈合六十点零五米，八丈五合十九点六米。}其二在材料：底部金属层的反射特性与月面样品中的一类合金十分接近。其三在位置：此类遗迹数量不多，分布范围大致与《左传》所记有穷氏的活动区域重合。这三条单拎出来都不算过硬，合在一处才勉强够格把地上与天上并作一件事。

一九七二年的那次测绘采用的是电阻率成像法：在地面布设九条测线，通以微弱电流，测量地下不同深度的电阻差异，空腔呈现高阻，金属呈现低阻，据此绘出地下轮廓，再结合反射波形确认底部存在金属层。\footnote{电阻率法只能给出异常的轮廓，无法分辨轮廓内的物质成分，因此底部「金属反射面」系结合反射波形共同判定。}

\begin{archive}[内部参考]{档案 № ФВ-1972-06}
地下空腔一处：深约十四米，长约六十米，宽约二十米，四壁规整。底部为金属反射面，当时完整，未见裂隙。
腔壁有刻痕，分四列，字不可识。
未取样，未进入。结论：不具备安全作业条件，建议维持封存。
\end{archive}

在这一处之外，尚有若干形制相近、尺寸不等的同类遗迹。\footnote{同类遗迹的清单与逐处尺寸见基金会勘察报告；本书只用得着有穷氏这一处，其余不赘。}勘察记录中有一句平淡无奇的话，却是本章的第二条物证：这些设施既无拆除痕迹，也无二次使用痕迹。没有拆除，说明他们动身时并不仓惶，一座准备彻底废弃的基地往往连管线都会拆走，这里却连管线出口都原样留着；没有二次使用，则说明他们一去再未复返，直到一九七二年测线布上地面，其间再无进人记录。至于他们为何离开、离开后又靠什么维生，地下并没有留下答案，腔壁上那四列刻痕至今也无人认得。

\section{月壳与船体}

离开地面的人去了何处，地下的遗迹没有明言，要接上这一环，还得回到第一章量过的那几片低地，并从整颗月球的物理参数来看。

先看低地本身。那几片低地的边界像是人为划出的，下方既未曾被撞击掘开，边缘又有垒筑的壁垒。在月球三号那二十九帧影像上，最大的几片低地之间以窄颈相连，最窄处约十二公里，窄颈两侧的壁垒依然完整，未见后期撞击打断的痕迹。几片互不相干的天然洼地大可不必有此窄颈，唯有连成一气的内部空间才需要通道串联。这一内部空间是否够住，第三章会细算：按城中环制与人均面积计，三十一万二千人所需地面以百平方公里计，而那几片低地连成一片，面积远在此数之上。

再看整颗月球，有三项反复测定的物理数值始终与「实心岩石天体」的说法格格不入。第一项是转动惯量：月球的转动惯量系数约为零点三九三，极接近均匀球体的零点四。一个充分分异的岩石天体，重元素理应沉降于核心，使其转动惯量系数明显低于零点四；数值如此贴近零点四，正说明它内部缺少一个致密的大核，若是内部中空、仅余外层厚壳，反而恰好能给出接近零点四的数值。\footnote{转动惯量系数由地月激光测距与月球轨道进动长期测定，精度已达千分位。数值偏近均匀球体的直接原因是核小，这一点学界一致；至于核为何如此之小则众说纷纭，本书取空壳说。}第二项是平均密度：月球平均密度约为每立方厘米三点三四克，而阿波罗与月球号带回的月岩样品密度多在三点一至三点五之间，玄武岩更是高于此数；浅层孔隙的贡献已被重力与地震波排除，剩下的密度缺口只能落在深部。第三项是月震：阿波罗任务在月球正面架设了四台地震仪，从一九六九年运行至一九七七年，记录下数千次月震；反演出的波速曲线并不平滑，不仅显示月球内部缺少一个像样的致密核心，还多出一层来历不明的低速带。\footnote{月震台网的四台地震仪均布设于月球正面，对背面的约束极弱，这是本节全部数据中最不确定的一处。}

三个数值分开看，各自都能找到别的地质解释；可合在一处，能同时容纳「几乎无核」「平均密度低于地表岩石」与「深部存在低速带」三件事的构造，大概只有一种，即一颗近乎空心的球体，外包一层厚壳。至于这层壳究竟多厚，本书给不出确数：月震反演给出的深度区间宽达一个数量级，背面钻探又只及地下四十米，眼下只能说它存在，说不清它的具体尺寸。

上述三项都是间接推论，更直接的证据来自《章程总目》的《居止》卷第一条。这条文字并非规程，倒像一段写实的记述。

\begin{archive}[内部参考]{档案 № ФВ-1965-09}
居止卷首条：「城在壳内。壳外无物，壳内无天。自上视之，壳为地；自内视之，地为壳。\\
壳厚若干，不载。\\
四时不由天，由机。」
\end{archive}

最初整理这批材料的人读不懂这一段，只当「壳」是指城中的某重环壁。按此读法，前两句还勉强说得过去，最后一句却说不通：城内四季由机械调节本是常识，何须再郑重写一句「四时不由天」。唯有把「壳」字读作月球本身，四句话才豁然贯通：城在月球内部，壳外别无他处可去；从外侧俯瞰，壳便是脚下的月面，从内部仰望，月壳便是头顶的穹天；末句「四时不由天，由机」，更是只有住在人造封闭空间里的人才写得出的句子。前三项物理参数说的只是「它像是一层壳」，这一条则直接点明「写章程的人就住在壳里」。

至此可以合拢为三层结构：月球是一层壳，壳内留有空腔，空腔中建有一座城。第一层由转动惯量、密度与月震给出，第二层由低地相连给出，第三层由章程自述给出。还差半步要判断这层壳究竟是天然形成还是人为建造。即便天然空腔勉强能解释无核与低速带，它也解释不了两件事：其一，一颗天然的空心天体何以恰好停在地球轨道上，且始终以同一面朝向地球；其二，低地的边界由人划定，环制的城由人修筑，天然洞穴不会自己长出一套环制城市。因此本书只能取一个看似荒诞的结论：月球并非天然卫星，而是一艘停泊在现今轨道上的船，外壳即今日所见的月面，一面朝向地球，背面朝外。

\section{灾难说与选择说}

面对同一批器物与岩芯，其实可以读出两种截然不同的历史。这里将两种读法一并列出，也各自指出它们对不上的地方。

一种读法是灾难说：地上突发变故，他们无法立足，于是造船远遁。这种读法顺理成章，最省力气，却有一处与岩芯对不上。BDP-96 孔六点零三米处的那一层泥并不算惨烈，植被退走之后数百年内硅藻便恢复如初，泥层里既无火成灰烬，也无洪水淤积或一次性生物灭绝带。若真是一场灭顶之灾，湖泥里理应留下更硬的灾难证据。

另一种读法是选择说：他们原本便有能力离去，且筹备已久，那场变故只是将动身的日子提前了。题头那一枚金属片支持了这一读法的一半：城池按三十一万二千人设计，登船名额却定在十五万，说明早在动身之前，留谁走谁的大数便已算定，并非临时清点。可这一读法同样有一处破绽：若早有周全准备，题头那三行字便不该写得如此仓促，连「存者如右」后方那一栏也留作空白。一份筹划多年的名单，照理不该空着最紧要的那一栏。

两种读法都无法严丝合缝地解释全部材料。笔者个人倾向选择说里的后半句，即「原本能走，只是提前走了」，理由仅此一条：它让地下遗迹「未见拆除痕迹」这一现象不必再费心找补，其余说法都得为那座完整封存的空腔另寻托词，而这一说不用。这一倾向并无实证支持，写在这里，只为让读者明白后面各章是踩着哪块石头往前走的。

\section{年代与推论的限度}

回过头来看，这一章摆上的无非三件东西：一根岩芯，一枚题头，一处地下遗迹，各自带着一个年代或一个姿态。岩芯显示距今约两万一千三百年曾有一场来势极快的植被更替；遗迹显示造访者走得从容，此后不曾折返；题头则显示他们登船前便算好了半数去留，只是未能留名。

三样合在一处，大概只能拼出一个轮廓：他们原居于此，约两万一千多年前离开，此后不再回来。至于离去的真正动因，现有材料里只字未提，本章给出的灾难与选择二说皆属猜测，且已标明各自的缺口。年代的尺度也必须如实交代：两万一千三百年的岩芯测年本身带着数百年误差，题头又全无纪年，两者跨历法换算尚有第七章要讲的几道关卡。把三件材料放在一起，眼下至多只能说彼此尚不抵牾，远谈不上彼此印证。

下一章将从那四百一十二枚金属片说起。那是这座城自己写下的文字，篇幅比题头的三行字大得多，也比题头更难读懂。
